\section{System}
\label{sec:system}

\subsection{Background}

GeoCLIP~\cite{cepeda2023geoclip} encodes an image $x$ with a CLIP ViT-L/14
vision tower $f_{\mathrm{img}}$ and GPS coordinates with a location encoder
$f_{\mathrm{loc}}$, then retrieves
\begin{equation*}
\hat{g}(x)=\arg\max_{g\in\mathcal{G}}
\langle\bar f_{\mathrm{img}}(x),\bar f_{\mathrm{loc}}(g)\rangle .
\end{equation*}
Here, $\bar{\cdot}$ denotes $L_2$ normalisation and $\mathcal{G}$ is a fixed
gallery of $10^{5}$ coordinates. GeoProbe leaves $f_{\mathrm{loc}}$ and
$\mathcal{G}$ unmodified throughout; every intervention acts upon
$f_{\mathrm{img}}$ alone, so that any change in $\hat{g}$ is attributable to
the image pathway.

\subsection{Region-Conditioned Attention Intervention}
\label{sec:intervention}

In each of the $L=24$ self-attention layers of the vision tower, attention
weights are obtained as $\mathrm{softmax}(QK^{\top}/\sqrt{d})$ over a
sequence of $257$ tokens: one \texttt{CLS} token followed by $256$ patch
tokens arranged on a $16 \times 16$ grid. Given a region mask
$m \in \{0,1\}^{256}$ indicating the patches falling within the selected
region, GeoProbe adds a key-indexed bias to the attention logits of layer
$\ell$ prior to the softmax,
\begin{equation}
  A^{(\ell)} = \mathrm{softmax}\!\left(
  \frac{QK^{\top}}{\sqrt{d}} + \beta^{(\ell)} \right),
  \qquad
  \beta^{(\ell)}_{j} =
  \begin{cases}
    0      & j = 0 \quad (\texttt{CLS}), \\
    a_\ell & j > 0,\; m_{j-1} = 1, \\
    b_\ell & j > 0,\; m_{j-1} = 0.
  \end{cases}
  \label{eq:bias}
\end{equation}

Three properties make this a suitable analytical primitive. The bias is
indexed by key and broadcast across all queries and heads, so every token
--- including \texttt{CLS} as a query --- draws proportionally more
information from within the region when $a_\ell > 0$ and less from outside
when $b_\ell < 0$. It is exactly a multiplicative reweighting of attention
mass by $\exp(\beta_j)$, so the control may be presented either as a logit
bias or as a scale factor. And $a_\ell = b_\ell = 0$ is an exact no-op, so
the reported baseline is the unmodified model rather than an approximation.
The \texttt{CLS} key is never biased: as a summary token rather than a
spatial patch, the interior/exterior distinction does not apply to it, and
biasing it would confound the regional effect with a global rescaling of the
summary pathway.
Here, ``causal'' refers narrowly to the effect of this controlled operation on
the model's computation; it does not by itself establish that a semantic object
causes the prediction in the photographed world.

\subsection{Mapping Regions to Patches}

A region delineated on the original image cannot be transferred to the patch
grid directly, since CLIP resizes the shorter side to $224$ pixels and then
applies a centre crop. GeoProbe replicates this geometry on the region
coordinates, so that a patch is marked in-region only if it remains so after
preprocessing. Regions falling entirely outside the centre crop yield an
empty mask and therefore an exact no-op. We retain these cases in the
full-dataset results and additionally report the active-mask subset. Multiple
regions are combined by union.

\subsection{Prompt-Free Region Proposals}

Manually delineated rectangles support targeted case studies, but they do
not scale to a complete dataset and make the analysis depend on a
user-supplied concept. For systematic evaluation, GeoProbe uses WeDetect-Uni,
a universal proposal generator trained with a learned objectness prompt to
retrieve category-agnostic foreground boxes~\cite{fu2026wedetect}. It
therefore avoids the per-image text query required by grounded detectors such
as Grounding DINO~\cite{liu2024grounding}. The detector
never observes ground-truth coordinates. Proposals are computed once and
cached; downstream layer searches only change the attention intervention, so
every configuration receives identical regions.

\subsection{Implementation}

A FastAPI backend loads GeoCLIP once at startup and patches the eager
attention forward pass of all $24$ vision layers according to
Eq.~\ref{eq:bias}; a React frontend provides region selection, per-layer
control and map display. One request returns both the unmodified and the
intervened top-$k$ predictions, so the comparison presented to the user is
always between two forward passes of identical weights. Manual rectangles
are used in the live interaction, while cached WeDetect-Uni proposals provide
the reproducible dataset-scale analysis reported below.
